\section{A centered height blow-up}\label{sec:centered-blowup}

The interval construction has supplied height comparison at the target
scales without changing the sequence that carries the density centers.
We now normalize height at those scales. The resulting centers have
mean curvature negligible at the height scale, and their first height
moments vanish. These two properties will allow both first variation
and the density comparison to pass to the limit.

\begin{proposition}[Centered blow-up at the target scales]
\label{prop:centered-blowup}
Let the stationary integral varifold, singular origin, integer $Q$,
plane $\Pi$, and sequence $L_i\downarrow0$ be as in
Proposition~\ref{prop:flat-scales}. After taking a subsequence and
identifying $\Pi$ with the horizontal plane, there are positive
$a_i\to0$ and $C^J$ graph disks $S_i$ over a fixed enlargement of
$B_8\subset\Pi$, converging to $\Pi$ through all derivative orders
needed below, with the following properties. In target units let
$\mu_i$ be the weight of $V_{L_i}$, let $\pi_i$ be nearest projection
to $S_i$, and set
\[
 z_i=X-\pi_iX,\qquad
 M_i=(\pi_i)_\#\mu_i-Q\vol_{S_i},\qquad
 D_i=(\pi_i)_\#(e_i\mu_i).
\]
All measures are restricted to one fixed interior normal cylinder
containing the support in question, with a support-free gap at its
vertical boundary. On the part over $B_8$,
\begin{equation}\label{eq:centered-interface}
 \begin{gathered}
 D_i+|M_i|=O(a_i^2)\quad\hbox{in total mass},
 \qquad |z_i|\le Ca_i\quad\hbox{on the support},\\
 \sup\bigl(|\mathcal H_{S_i}|+
                  |\nabla^\perp\mathcal H_{S_i}|\bigr)=o(a_i),
 \qquad
 a_i^{-1}(\pi_i)_\#(z_i\mu_i)\rightharpoonup0.
 \end{gathered}
\end{equation}
Moreover the normalized height-square measures have uniformly bounded
mass there and are quantitatively nonzero on a fixed compact subset
of $B_4$; in particular, for all sufficiently large $i$,
\begin{equation}\label{eq:nonzero-second-moment}
 c\le a_i^{-2}\int_{\pi_iX\in B_3}|z_i|^2\,\dd\mu_i
 \le a_i^{-2}\int_{\pi_iX\in B_8}|z_i|^2\,\dd\mu_i\le C.
\end{equation}
Balls in this statement are understood through the graph coordinates
over $\Pi$. The weak convergence in \eqref{eq:centered-interface}
holds against compactly supported tests in $B_8$.
\end{proposition}

\begin{proof}
We prove successively the integral normalization, negligible curvature,
uniform support height, and first-moment centering. In all intermediate
estimates we retain larger fixed interior regions than those in the
statement.

\subsubsection*{Integral normalization and convergence of the centers}
For the interval covering $L_i$, let $R$ be its start, use its center
$S$, and put
\begin{equation}\label{eq:height-normalization}
 r_i=L_i/R,\qquad a_i=\frac{\sqrt{H(r_i)}}{r_i}>0.
\end{equation}
The notation $R$ and $S$ here may depend on $i$. If $r_i\to0$,
then $a_i\to0$ by \eqref{eq:height-power-bounds} and
$\delta\le Cr_i^{k-1}$. If $r_i$ stays bounded away from zero,
we are in the multiple-interval case, and
\eqref{eq:restart-height-upper}, with $h_0\to0$, gives the same
conclusion. These cases also show that $\delta(cr_i)\to0$ for every
fixed enlargement used below. Subsequences suffice to handle them.

Proposition~\ref{prop:target-frequency} yields
\begin{equation}\label{eq:target-energy}
 \begin{aligned}
 r_i^{-m}\int_{\rho<200r_i}|z|^2\,\dd\mu
        &\le Ca_i^2r_i^2,\\
 r_i^{-m}D(B^S_{200r_i})&\le Ca_i^2,
 \end{aligned}
\end{equation}
where these expressions remain in start units. Indeed
\eqref{eq:height-doubling} supplies a radius in
$[400r_i,600r_i]$ at which $q_*$, and hence $n_*$, is bounded.
These radii lie below $r_*$ because $r_i\le r_*/1000$.
The weights at that radius are bounded below on
$\{\rho<200r_i\}$, so
\eqref{eq:weighted-height-comparisons} gives the height bound, and
the definition of $E_*$ gives the tilt bound. There is also a converse
height lower bound $ca_i^2r_i^2$ already on $\{\rho<2r_i\}$,
directly from $H(r_i)=a_i^2r_i^2$.

Dilate the center into target units and denote it by $S_i$.
Its graphs have uniform high derivative bounds on fixed balls: their
start derivatives are uniformly bounded and the dilation introduces
factors $r_i^{j-1}$ in derivative order $j$. Their heights at
horizontal $0$ before dilation are at most $Cm_0b^k$, and their slopes
are small. Since $b\le r_i/B_*$, the center positions remain bounded
after dilation. The sup height estimates show that support points on
each needed fixed region are at distance tending to zero from these
graphs in target units. For example, support points in an ambient
ball of radius $60$ have projected geodesic radius below $180$,
by the fixed projection bounds and the position of the radial center.
Support convergence to $Q|\Pi|$ therefore places a corresponding
patch of $\Pi$ in every subsequential graph limit. A small-slope graph
containing this patch must equal the plane there. Compactness with
extra derivative orders gives the asserted smooth-enough convergence.
Passing to a subsequence also fixes the limiting rotations of the start
coordinates. We may now regraph over $\Pi$ on smaller fixed room.

\subsubsection*{Curvature below the height scale}
We claim that, on a fixed region larger than the one in the statement,
\begin{equation}\label{eq:center-curvature}
 \sup\bigl(|\mathcal H_{S_i}|+
                 |\nabla^\perp\mathcal H_{S_i}|\bigr)=o(a_i).
\end{equation}
Temporarily return to start units. At a point with horizontal argument
of norm at most $40r_i$ and $t>0$, choose a witness $(x,\ell)$.
The argument for \eqref{eq:witness-size} gives
$\ell\le\max(b,40r_i/15)\le3r_i$.
Apply failed-fit separation on a small subball near that point; its
center lies within the allowed $30\ell$ of $x$. The subball, and
the changes between its fitting and normal projections, remain inside
the upper-bound region in \eqref{eq:target-energy}. Separation and
small slopes give height at least $cm_0\ell^k$ on mass at least
$c\ell^m$. Consequently
\begin{equation}\label{eq:witness-normalization}
 a_i r_i\ge cm_0\ell^k(\ell/r_i)^{m/2}.
\end{equation}
All such witness scales tend to zero uniformly. This is immediate
when $r_i\to0$, and otherwise follows from $a_i\to0$ in
\eqref{eq:witness-normalization}.

The pointwise start-unit curvature estimate is
$Cm_0^2\ell^{2k-4}$, by \eqref{eq:glued-center} and $t\le2\ell$.
Mean curvature and its first normal derivative acquire factors
$r_i$ and $r_i^2$, respectively, in target units. Thus, for mean
curvature, \eqref{eq:witness-normalization} gives
\[
 \frac{r_i m_0^2\ell^{2k-4}}{a_i}
       \le Cm_0r_i^{2+m/2}\ell^{k-4-m/2}\longrightarrow0.
\]
The derivative has one additional factor $r_i\le1$. The exponent of
$\ell$ is positive by \eqref{eq:fitting-order}. At $t=0$ both
curvature quantities vanish by \eqref{eq:glued-center}. This proves
\eqref{eq:center-curvature}.

\subsubsection*{A uniform bound for support height}
Integral control alone would not suffice for all later expansions.
We strengthen it to
\begin{equation}\label{eq:target-height}
 |z_i|\le Ca_i
\end{equation}
on fixed inner room. The normal-displacement trace formula
\cite[Equation~(7.9)]{AE} is
\begin{equation}\label{eq:normal-displacement-trace}
 P_T:Dz_i=e_i-\mathcal H_{S_i}\cdot z_i
       +O(|\II_{S_i}|^2|z_i|^2+|\II_{S_i}|e_i|z_i|).
\end{equation}
Since $\nabla|z_i|^2=2z_i$, the nonnegative ambient function
$G=|z_i|^2+a_i^2$ satisfies, on the support,
\begin{equation}\label{eq:height-subsolution}
 \tr(P_TD^2G)\ge-CG.
\end{equation}
Here the positive $2e_i$ term absorbs its small multiple in
\eqref{eq:normal-displacement-trace}, and
\eqref{eq:center-curvature} bounds the term linear in height by
$C(|z_i|^2+a_i^2)$.

For completeness, \eqref{eq:height-subsolution} has the requisite
mean-value consequence on a stationary varifold. Fix a support point
$X_*$, put $d=|X-X_*|$, and take a nonnegative decreasing smooth
radial kernel $f(d/l)$, supported in $d<l$ and positive at $d=0$.
Define
\[
 I(l)=l^{-m}\int f(d/l)G\,\dd\mu_i.
\]
Radial stationarity gives
\[
 \frac{\dd I}{\dd\log l}
 \ge l^{-m}\int f(d/l)\nabla_TG\cdot(X-X_*)\,\dd\mu_i.
\]
The discarded term is nonnegative because $f$ is decreasing.
Use the primitive
$F_l(d)=\int_d^\infty t f(t/l)\,\dd t$, whose gradient is
$-f(d/l)(X-X_*)$ and which satisfies
$F_l\le Cl^2f(d/l)$. Stationarity and
\eqref{eq:height-subsolution} bound the right-hand side below by
$-Cl^2I(l)$. Integration of this differential inequality down to zero, using
density at least one, bounds $G(X_*)$ by a fixed-radius integral
of $G$ around $X_*$. The integral is $O(a_i^2)$ by
\eqref{eq:target-energy} and the mass bound, proving
\eqref{eq:target-height}.

These tests may be made for all points with horizontal arguments in a
fixed ball of radius $25$ in target units, using small fixed ambient
balls. Their heights already tend to zero by the earlier sup height
estimate. To check that the normal coordinates include all relevant
support, in start units use $t(w)\le b+|w|$ and
\eqref{eq:glued-center}: on these arguments the heights are small
compared with $r_i$, and changes from horizontal to projected arguments
are smaller still by the small slope. Decreasing $r_*$ once fixes all
needed tube and chart margins. Regraphing over $\Pi$ preserves smaller
fixed interior regions.

\subsubsection*{Vanishing of the first height moment}
It remains to prove
\begin{equation}\label{eq:centered-first-moment}
 a_i^{-1}(\pi_i)_\#(z_i\mu_i)\rightharpoonup0.
\end{equation}
We first use horizontal tests and vertical graph-height differences in
the start coordinates. Take an arbitrary fixed smooth test of $w/r_i$,
supported in horizontal $B_{12r_i}$, so that its derivatives are bounded
at that scale. We do not yet pull tests back through the changing
centers. By convergence of their coordinate changes, this region
contains the supports of target tests on $B_8\subset\Pi$ for all
large $i$; coordinate conversion will be made after proving the
fixed-coordinate limit.
Cover $\{t>0\}$ there by balls of radii a small fixed multiple of
$t$ at their centers, selecting a disjoint smaller family and a
bounded-overlap enlargement. The Lipschitz property of $t$ gives a
partition of unity whose derivatives are bounded at the corresponding
cell scales. Height is zero off $\{t>0\}$.

On a cell choose a constructed fit of dyadic scale $L$ comparable to
$t$ and sufficiently large compared with it, as in the gluing
construction. Cells meeting the test support have centers in an
$O(r_i)$ region. More explicitly, choose the enlarged cell radius
at most $ct(w)$ with $c<1/10$. A cell meeting $B_{12r_i}$ has
\[
 |w|\le12r_i+ct(w)\le12r_i+c(b+|w|),
 \qquad |w|<14r_i,
\]
because $b\le r_i/100$. In particular $L\le Cr_i$ and these centers
lie strictly inside the $40r_i$ region used for
\eqref{eq:witness-normalization}. Applying the
first-moment estimate in Lemma~\ref{lem:local-fits} to the partitioned
test, whose required scaled derivatives are bounded, gives an error
at most
\begin{equation}\label{eq:cell-moment-error}
 CL^m m_0^2L^{2k-1}.
\end{equation}
This controls the height over the local fit. Its difference from the
glued center has height $Cm_0^2L^{2k-1}$, so the same bound controls
that remaining contribution, using the cell mass bound.

To compare \eqref{eq:cell-moment-error} with the target height scale,
choose a witness at the cell center. It satisfies $\ell\ge t/2$,
hence $L\le C\ell$, though no converse comparison is required.
Equation~\eqref{eq:witness-normalization} implies
\[
 \frac{m_0^2L^{2k-1}}{a_i r_i}
    \le Cm_0r_i^{m/2}\ell^{k-1-m/2}\longrightarrow0
\]
uniformly over the cells. Positivity of the exponent follows from
\eqref{eq:fitting-order}, and the witness scales tend uniformly to
zero as proved above. Independently, bounded overlap of the cells
and $L\asymp t$ give $\sum L^m\le Cr_i^m$. Summing
\eqref{eq:cell-moment-error} therefore gives an error
$o(a_i r_i^{m+1})$. This is exactly the vanishing first moment after
mass and height are expressed in target units.

Finally convert from vertical graph-height differences to normal
displacements. At first order this conversion is the normal projection
at the graph point, with uniformly bounded quadratic error. By
\eqref{eq:target-height} the vertical differences also have size
$O(a_i)$. The quadratic error is therefore $O(a_i^2)$, and test
arguments change by $O(a_i)$, which has the same harmless effect.
The coefficient matrices converge uniformly on the fixed chart.
The normalized first-moment measures have bounded total variation,
so their convergence against fixed smooth tests extends to these
converging tests and matrices. This proves
\eqref{eq:centered-first-moment}.

We can now collect the conclusions in target units. The tilt estimate
in \eqref{eq:target-energy} gives $D_i=O(a_i^2)$ on the required
room. The scalar estimate, Lemma~\ref{lem:scalar-input}, with
\eqref{eq:center-curvature} and \eqref{eq:target-height}, gives
$|M_i|=O(a_i^2)$ there as well. Its integer is $Q$, identified by
smooth averages and convergence to $Q|\Pi|$. If a fixed small physical
radius is needed for that estimate, a fixed dilation changes only
constants. The lower height bound after
\eqref{eq:target-energy}, originally in geodesic radius $2r_i$,
lies in graph-coordinate $B_3$ for all large $i$, because the centers
converge to $\Pi$ and their radial origins tend to $0$.
This proves \eqref{eq:nonzero-second-moment}. All calculations retained
enough room to select the fixed cylinder over $B_8$, with a vertical
cutoff in a support-free gap. Hence they concern the full local
varifold measure required in the statement, and the proposition follows.
\end{proof}
